\section{Psychology of Worldviews}

There are two kinds of people in the world:

\begin{itemize}


\item those who try to persuade others

\item and those who live their lives as though nothing is going on beyond their personal concerns.
\end{itemize}


As for the first kind, it should at some point dawn on them that persuasion is singularly ineffective. They may try to persuade by reason or logic, but this has no effect on those who are influenced by other factors, such as emotion, group loyalties, and so on. A man's identity is intricately tied up with his worldview, so he is loathe to give it up. There are very few who attempt the First Trial\footnote{\url{https://www.gornahoor.net/?p=1902}}. Even the persuader is limited by his own worldview. This is because there are many personality types and those of different types are opaque to each other. They respond to different centers and have different orientations to the world.


\paragraph{Centers}


In his interior life, a man may be centered in the:

\begin{itemize}
\item \textbf{Head}: His consciousness is dominated by thinking.


\item \textbf{Heart}: His consciousness is dominated by feeling.


\item \textbf{Belly}: His consciousness is dominated by willing.

\end{itemize}


\paragraph{Orientation}


In his relation to the world, a man may have three orientations:

\begin{itemize}
\item \textbf{Extraversion}: He is oriented to the world outside the self.

\item \textbf{Intraversion}: He is oriented to the world within the self.

\item \textbf{Supraversion}: He is oriented to the absolute, beyond the self.
\end{itemize}


\paragraph{Objectivity}


That yields nine fundamental types, and many more subtypes, since a man is rarely in a single classification all the time. These nine types correspond to nine possible areas of interest, although it is difficult to get a man to think beyond his primary and secondary areas of interest. Another factor is whether he can experience things objectively, as they are in themselves, or subjectively, insofar as they affect him. This results in three more classifications:

\begin{itemize}


\item \textbf{Head}: He may be interested in personal opinions rather than what is True.

\item \textbf{Heart}: He may be interested in pleasure rather than Beauty.

\item \textbf{Belly}: He may be interested in what is useful rather than what is Good.
\end{itemize}


\paragraph{Conclusion}


\textbf{Alfred North Whitehead} wrote that a man wants to live, then live well, then live better. Hence, his goals are prosperity, justice in his relationships, and then eternal bliss. A man has difficulty in balancing those goals. For those focused on prosperity, for example, the higher goals are usually invisible to him. The spiritual path, therefore, will differ from person to person. As \textbf{William Blake} wrote: ``One law for the lion and ox is oppression.''

A young man asked me, earlier today, for advice on his own spiritual path. His mind was filled with conflicting demands from parents, peers, and so on, thus his own dharma was hidden from him. Of course, I couldn't tell him what to do. As I explained to him: ``It would be unwise of me to offer such advice to a new acquaintance and imprudent of you to accept it from a stranger on the Internet.'' Instead, I encouraged him to clear his mind of those nagging opinions in his head to let his own Will guide him, and suggested some spiritual exercises. Obviously, to determine one's True Will is difficult in the best of times, and all the more so in the modern world that does not have a place for those of a Traditional cast of mind.

A good spiritual exercise, then, is to try to determine the personality type of those you interact with during the day.


\flrightit{Posted on 2011-12-10 by Cologero}

\begin{center}* * *\end{center}

\begin{footnotesize}
\begin{sffamily}
\texttt{Timotheos on 2016-05-23 at 15:07 said:}

This is interesting but I don't see how to determine the personality types of those I interact with from this? Are you able to offer any further guidance please?

\hfill

\texttt{Cologero on 2016-05-23 at 20:08 said:}

Timotheos, this is one of the topics we discuss on Monday nights.

\hfill

\texttt{Timotheos on 2016-05-24 at 04:26 said:}

Pardon me but where do you discuss that?

\hfill

\texttt{francismercuri on 2016-05-24 at 22:14 said:}

Off the cuff, an interesting parallel exists within exoteric philosophy. For Epictetus, how one acts (and attempts to order their interior world) best reflects which ``school of philosophy'' one belongs too (esoterically, this is transformed into which ``centers'' are dominant). Epictetus explained that if one ``observes'' themselves daily, they will quickly discover their native/inherent ``school''--most of his students he said were still quite Epicurean in practice (pleasure=happiness); while others were better along, being internally Peripatetic (virtue is the good, but requiring that outside conditions are favorable to the individual in his pursuit). Quite dismally perhaps--yet realistically, he concluded that few if any of his students were really of the orientation (in practice) that Arete/Virtue alone is sufficient.

\hfill

\texttt{Matt on 2018-12-10 at 14:58 said:}

Being oriented to the absolute (supraversion) seems identical with experiencing things objectively. If so, the ability to experience things objectively can probably be collapsed into the orientation of supraversion. Or perhaps one could move supraversion into objective experiencing. Either way, I cannot imagine supraversion coexisting with experiencing things subjectively.

But I like the addition of supraversion to introversion and extraversion, which are normally characterized as two ends of a polar spectrum. Supraversion would be the point above and in the middle, making a triangle configuration. Thus, instead of being limited to a horizontal range of subjective orientation preference (which we are told we can't really change anyway), one can orient oneself objectively, toward the absolute. Something which I believe can be freely chosen.


\hfill

\end{sffamily}
\end{footnotesize}

